\documentclass[11pt,oneside]{book}
\usepackage{uwnotes}

\title{The Mathematics of the Music Microscope}
\author{Harry Geng}
\date{}

\begin{document}

% ================= title page =================
\begin{titlepage}
\centering
{\Huge\bfseries The Mathematics of the\\[4pt] Music Microscope\par}
\vspace{1.2cm}
{\Large Course Notes for the Audio-Visualizer Project\par}
\vspace{2.2cm}
{\large Harry Geng\par}
\vfill
{\small Living document --- grows with the codebase.\\
Chapters marked \emph{(outline)} are planned but not yet written.\par}
\vspace{0.8cm}
{\small \today\par}
\end{titlepage}

\frontmatter
\tableofcontents

% ================= preface =================
\chapter{Preface}
These notes document, with proofs where they are illuminating, the mathematics
running inside the Audio-Visualizer project: a system that ingests songs,
separates them into stems with neural networks, extracts musical structure,
embeds every few seconds of audio into learned vector spaces, and renders the
result as interactive visualizations.

The material spans four traditions that rarely share a book: classical signal
processing (Part~\ref{part:signals}), modern deep learning
(Part~\ref{part:learning}), the geometry of high-dimensional data
(Part~\ref{part:geometry}), and computer graphics
(Part~\ref{part:rendering}). The unifying thread is that \emph{every object
defined here actually runs}: each chapter ends where the repository begins,
and the ``In the code'' boxes point at the specific file, function, and
parameter values where a definition earns its keep. When a theorem's
hypothesis is violated in practice (they often are), the notes say so and
describe what breaks.

Exercises range from short proofs to experiments you can run against your own
music library. None are busywork; several reproduce bugs that were actually
encountered while building the system.

\section*{How to read this}
Chapter~\ref{ch:sound} and Chapter~\ref{ch:timbre} are complete and set the
standard for the rest. The remaining chapters exist as detailed outlines and
will be filled in as the corresponding subsystems stabilize. Prerequisites are
first-year linear algebra and calculus; everything else is developed as
needed.

\section*{Notation}
\begin{center}
\begin{tabular}{ll}
\toprule
$x[n]$ & a discrete-time signal (square brackets: integer index) \\
$f_s$ & sampling rate in Hz \\
$X[k]$ & the DFT of $x[n]$; $k$ indexes frequency bins \\
$\norm{x}$ & Euclidean norm; $\norm{x}_F$ the Frobenius norm \\
$\ip{x,y}$ & inner product $x^\T y$ \\
$S^{d-1}$ & the unit sphere $\{x \in \R^d : \norm{x} = 1\}$ \\
$A^\T$ & matrix transpose \\
$\mathbf{1}[\cdot]$ & indicator function \\
$\log$ & natural logarithm unless a base is written \\
\bottomrule
\end{tabular}
\end{center}

\mainmatter

% ================= PART I =================
\part{Signals}\label{part:signals}
\chapter{From Sound to Vectors}\label{ch:sound}
\chapterlede{Everything in this project begins by turning pressure waves into
arrays of numbers, and arrays of numbers into frequencies. This chapter builds
the discrete Fourier transform from scratch, proves the identities the rest of
the book leans on, and ends at the two workhorse measurements the visualizer
makes sixty times per second: the spectrogram and the autocorrelation pitch
estimate.}

\section{Sampling}

A microphone measures a continuous function $x_c : \R \to \R$ (air pressure
against time). Computers store countably many numbers, so we sample.

\begin{definition}[Sampling]\label{def:sampling}
Given a continuous signal $x_c$ and a \emph{sampling rate} $f_s > 0$ (in Hz),
the sampled signal is the sequence
\[
x[n] = x_c\!\left(\tfrac{n}{f_s}\right), \qquad n \in \Z .
\]
\end{definition}

Sampling destroys information in general: infinitely many continuous signals
agree on any given sample sequence. The classical theorem says that
\emph{bandlimited} signals are the exception.

\begin{theorem}[Nyquist--Shannon]\label{thm:nyquist}
Let $x_c$ be bandlimited to $B$ Hz, meaning its continuous Fourier transform
vanishes outside $[-B, B]$. If $f_s > 2B$, then $x_c$ is uniquely determined
by its samples $x[n]$, and can be reconstructed as
\[
x_c(t) \;=\; \sum_{n \in \Z} x[n] \,
\operatorname{sinc}\!\big(f_s t - n\big),
\qquad \operatorname{sinc}(u) = \frac{\sin(\pi u)}{\pi u}.
\]
\end{theorem}

\begin{proof}[Proof idea]
Bandlimitedness means the Fourier transform $\hat{x}_c$ is supported on
$[-B,B]$. Sampling in time corresponds to \emph{periodizing} in frequency:
the transform of the sample train is $\sum_k \hat{x}_c(f - k f_s)$, a sum of
copies of $\hat{x}_c$ spaced $f_s$ apart. If $f_s > 2B$ the copies do not
overlap, so multiplying by the indicator of $[-f_s/2, f_s/2]$ recovers
$\hat{x}_c$ exactly; that multiplication in frequency is convolution with the
sinc kernel in time.
\end{proof}

When $f_s \le 2B$ the copies overlap and distinct frequencies become
indistinguishable after sampling --- \emph{aliasing}: a tone at frequency $f$
is captured identically to one at $f_s - f$.

\begin{incode}
The project stores two versions of every stem
(\texttt{ingest.py}): a listening copy at $f_s = 44{,}100$ Hz stereo
(Nyquist frequency $22.05$ kHz, above human hearing) and an analysis copy at
$f_s = 22{,}050$ Hz mono (\texttt{config.py}: \texttt{SR = 22050}), which is
cheaper to process and still covers content up to $11$ kHz --- plenty for
beat, pitch, and chroma analysis. Every resample (e.g.\ to MERT's $24$ kHz in
\texttt{embeddings.py}, or CLAP's $48$ kHz in \texttt{backfill\_clap.py})
must low-pass first for exactly the reason in
Theorem~\ref{thm:nyquist}; \texttt{librosa.resample} does this internally.
\end{incode}

\section{The Discrete Fourier Transform}

Fix a frame length $N$ and view a chunk of signal as a vector
$x = (x[0], \dots, x[N-1]) \in \C^N$.

\begin{definition}[DFT]\label{def:dft}
Let $\omega_N = e^{-2\pi i / N}$. The \emph{discrete Fourier transform} of
$x \in \C^N$ is $X \in \C^N$ with
\[
X[k] \;=\; \sum_{n=0}^{N-1} x[n]\, \omega_N^{kn}
\;=\; \sum_{n=0}^{N-1} x[n]\, e^{-2\pi i k n / N},
\qquad k = 0, \dots, N-1 .
\]
\end{definition}

The DFT is the change of basis into the family of sampled complex sinusoids
$e_k[n] = e^{2\pi i k n / N}$. That family is orthogonal:

\begin{lemma}[Orthogonality]\label{lem:orth}
For $k, \ell \in \{0, \dots, N-1\}$,
\[
\sum_{n=0}^{N-1} e^{2\pi i (k-\ell) n / N} =
\begin{cases}
N & k = \ell,\\
0 & k \ne \ell.
\end{cases}
\]
\end{lemma}

\begin{proof}
If $k = \ell$ every term is $1$. Otherwise let
$q = e^{2\pi i (k-\ell)/N} \ne 1$; the sum is geometric:
$\sum_{n=0}^{N-1} q^n = \frac{q^N - 1}{q - 1}$,
and $q^N = e^{2\pi i (k - \ell)} = 1$, so the numerator vanishes.
\end{proof}

\begin{corollary}[Inversion and Parseval]\label{cor:parseval}
$x[n] = \frac{1}{N}\sum_{k=0}^{N-1} X[k]\, e^{2\pi i k n/N}$, and
\[
\sum_{n=0}^{N-1} \abs{x[n]}^2 \;=\; \frac{1}{N} \sum_{k=0}^{N-1}
\abs{X[k]}^2 .
\]
\end{corollary}

\begin{proof}
Substitute Definition~\ref{def:dft} into the claimed inverse and swap the two
finite sums; Lemma~\ref{lem:orth} collapses the inner sum to $N\,
\mathbf{1}[n = m]$. Parseval follows by expanding
$\abs{X[k]}^2 = X[k]\overline{X[k]}$ and applying the lemma once more.
\end{proof}

Parseval is the statement that \emph{energy is preserved}: what the ear hears
as loudness can be measured either per-sample or per-frequency-bin. The
visualizer uses both sides of this identity (RMS level on the time side, band
energy on the frequency side).

\begin{remark}[What a bin means]\label{rem:bins}
Bin $k$ of an $N$-point DFT at sampling rate $f_s$ represents frequency
\[
f_k = \frac{k \, f_s}{N} \quad \text{Hz}, \qquad k \le N/2,
\]
so the frequency resolution is $\Delta f = f_s / N$: longer frames see finer
pitch differences. For a real signal, $X[N-k] = \overline{X[k]}$, so only the
first $N/2 + 1$ bins carry information.
\end{remark}

\begin{incode}
The chroma computation in \texttt{scene.js}
(\texttt{computeChroma}) inverts Remark~\ref{rem:bins}: it walks the
analyser's bins and maps each to its frequency via
\texttt{fr = bin * sr / nfft} before folding onto pitch classes
(Appendix~\ref{ch:temperament}). With \texttt{fftSize = 1024} at
$f_s = 44.1$ kHz, $\Delta f \approx 43$ Hz --- coarse at the bottom of the
piano (the gap between $A1$ and $B\flat1$ is only $3.3$ Hz) and luxurious at
the top. This asymmetry between linear bins and logarithmic musical pitch
recurs throughout the project.
\end{incode}

\section{The Fast Fourier Transform}

Computing Definition~\ref{def:dft} literally costs $N$ multiplications for
each of $N$ bins: $\Theta(N^2)$. The FFT gets the same numbers in
$\Theta(N \log N)$.

\begin{theorem}[Cooley--Tukey]\label{thm:fft}
Let $N$ be even, and split $x$ into even- and odd-indexed halves with DFTs
$E$ and $O$ (each of length $N/2$). Then for $k = 0, \dots, N/2 - 1$:
\[
X[k] = E[k] + \omega_N^{k}\, O[k], \qquad
X[k + N/2] = E[k] - \omega_N^{k}\, O[k].
\]
Consequently, for $N = 2^m$, the DFT can be computed in
$\Theta(N \log N)$ operations.
\end{theorem}

\begin{proof}
Split the defining sum by parity of $n$, writing $n = 2r$ and $n = 2r + 1$:
\[
X[k] = \sum_{r=0}^{N/2-1} x[2r]\, \omega_N^{2rk}
     + \omega_N^{k} \sum_{r=0}^{N/2-1} x[2r+1]\, \omega_N^{2rk}
     = E[k] + \omega_N^k O[k],
\]
using $\omega_N^2 = \omega_{N/2}$. The second identity follows from
$\omega_N^{k + N/2} = -\omega_N^{k}$. The cost recurrence
$T(N) = 2\,T(N/2) + \Theta(N)$ solves to $\Theta(N \log N)$ by the master
theorem.
\end{proof}

This is why every frame size in the project is a power of two
(\texttt{fftSize} of $1024$, $2048$): the browser's
\texttt{AnalyserNode} and NumPy's FFT both rely on this factorization.

\section{Frames, Windows, and the Spectrogram}

Music is not stationary --- the whole point is that its frequency content
changes. The fix is to transform short \emph{frames} and accept a tradeoff.

\begin{definition}[STFT and spectrogram]\label{def:stft}
Fix a frame length $N$, a hop $H \le N$, and a window
$w : \{0,\dots,N-1\} \to [0,1]$. The \emph{short-time Fourier transform} of
$x$ is
\[
X[m, k] = \sum_{n=0}^{N-1} x[mH + n]\, w[n]\, e^{-2\pi i k n / N},
\]
and the \emph{spectrogram} is the matrix of squared magnitudes
$S[m,k] = \abs{X[m,k]}^2$.
\end{definition}

Without a window ($w \equiv 1$), chopping the signal introduces artificial
discontinuities at frame edges, and Definition~\ref{def:dft} smears each true
frequency across all bins (\emph{spectral leakage}). A smooth taper such as
the Hann window
\[
w[n] = \tfrac{1}{2}\Big(1 - \cos\tfrac{2\pi n}{N-1}\Big)
\]
concentrates that leakage into nearby bins at the price of slightly widening
each spectral peak.

\begin{remark}[Time--frequency tradeoff]\label{rem:uncertainty}
$\Delta f = f_s / N$ improves with longer frames while time localization
$\Delta t = N / f_s$ worsens \emph{in exact inverse proportion}:
$\Delta t \cdot \Delta f = 1$ regardless of $N$. This is the discrete shadow
of the uncertainty principle. There is no window that beats it; one can only
choose where to spend it. Musical consequence: a system tuned to see bass
notes precisely (long frames) necessarily blurs drum onsets, and vice versa.
\end{remark}

\begin{incode}
The microscope's spectrogram view (\texttt{microscope.py},
\texttt{spectrogram\_png}) renders $\log(1 + S)$ rather than $S$ --- the log
compresses the enormous dynamic range of music ($\sim$60 dB between a hi-hat
tail and a kick) into the few dozen brightness levels an eye can
distinguish, then maps it through the magma colormap. The scene mode's live
analysis instead uses the browser's \texttt{AnalyserNode}
(\texttt{app.js}, \texttt{bands()}), which computes
Definition~\ref{def:stft} with a built-in Blackman window at
\texttt{fftSize = 1024}: $\Delta t \approx 23$ ms, $\Delta f \approx 43$ Hz
--- the ``spend it on time'' end of Remark~\ref{rem:uncertainty}, correct for
a visualizer chasing drum hits.
\end{incode}

\section{Pitch by Autocorrelation}

The spectrogram spreads a note's energy across its harmonics; for
\emph{monophonic} sources (one note at a time --- bass, a single voice) a
time-domain method finds the fundamental more directly.

\begin{definition}[Autocorrelation]\label{def:acf}
The (unnormalized, windowed) autocorrelation of $x$ at lag $\tau \ge 0$ over
a window of length $W$ is
\[
r[\tau] = \sum_{n=0}^{W-1} x[n]\, x[n + \tau].
\]
\end{definition}

\begin{proposition}[Periodicity peaks]\label{prop:acfpeak}
Suppose $x$ is $P$-periodic on the relevant range, i.e.\ $x[n + P] = x[n]$.
Then $r[\tau] \le r[0]$ for all $\tau$, with equality at every multiple of
$P$: $r[0] = r[P] = r[2P] = \cdots$
\end{proposition}

\begin{proof}
Equality at multiples of $P$ is immediate from periodicity. For the bound,
Cauchy--Schwarz gives
\[
r[\tau] = \ip*{\,x_{0:W},\, x_{\tau:\tau+W}}
\le \norm{x_{0:W}} \, \norm{x_{\tau:\tau+W}} = r[0],
\]
where the final equality uses periodicity again (both windows contain the
same multiset of values when $W$ is a whole number of periods).
\end{proof}

So the algorithm is: compute $r[\tau]$ over a plausible range of lags, find
the peak, and report $f_0 = f_s / \tau^\ast$. Two practical failure modes
matter. First, $r$ peaks equally at $2P, 3P, \dots$, so noise can make the
estimator jump down an octave (choosing $2P$) --- the classic \emph{octave
error}. Second, on silence or noise the global maximum is meaningless, so an
energy gate must precede the search.

\begin{incode}
\texttt{Transport.pitchHz} in \texttt{app.js} is
Definition~\ref{def:acf} verbatim: window $W = 512$ samples, lag range
$f_s/320 \le \tau \le f_s/50$ (i.e.\ candidate fundamentals between $50$ and
$320$ Hz --- the playable range of a bass guitar), and an RMS gate of
$0.012$ before searching. The stutter you once saw in the bass fret marker
was octave error plus frame-to-frame noise; the mitigations in
\texttt{scene.js} (sampling every third frame, level-gating, exponential
smoothing $0.7/0.3$) treat the symptoms. The planned cure --- an offline
pitch tracker with Viterbi smoothing across the whole song --- replaces the
per-frame $\argmax$ with a globally optimal path and will live in the ingest
pipeline.
\end{incode}

\section*{Exercises}

\begin{exercise}
A pure $9$ kHz tone is sampled at $f_s = 16$ kHz with no anti-aliasing
filter. What frequency is heard on playback? Generalize: give the perceived
frequency as a function of true frequency $f$ and $f_s$.
\end{exercise}

\begin{exercise}
Prove the conjugate-symmetry claim of Remark~\ref{rem:bins}: if
$x[n] \in \R$ for all $n$, then $X[N - k] = \overline{X[k]}$. How many real
numbers are needed to store the DFT of a real signal, and why does this not
contradict the DFT being invertible?
\end{exercise}

\begin{exercise}
Complete the proof of Corollary~\ref{cor:parseval} in full detail, justifying
the interchange of summations.
\end{exercise}

\begin{exercise}
Show that the Hann window satisfies
$\sum_{m \in \Z} w[n - mH] = \text{const}$ when $H = N/2$ (perfect overlap-add
at 50\% hop). Why does this property matter for reconstructing audio from a
modified STFT --- for instance, in mask-based source separation
(Chapter~\ref{ch:separation})?
\end{exercise}

\begin{exercise}[computational]
Record or load a bass stem from your library. Implement
Definition~\ref{def:acf} and plot $r[\tau]$ for one frame during a held note.
Identify the fundamental's peak and at least one octave-error peak. Then
break the estimator on purpose: find a frame where the $\argmax$ lands at
$2P$.
\end{exercise}

\begin{exercise}
Using Remark~\ref{rem:uncertainty}: the microscope wants to distinguish
$E1$ ($41.2$ Hz) from $F1$ ($43.7$ Hz) in the bass register. What minimum
frame length (in samples at $22{,}050$ Hz, and in milliseconds) does a plain
DFT need? Roughly how many kick drum hits at $170$ BPM occur within one such
frame? Conclude something about why the project analyzes bass pitch and drum
onsets with different tools.
\end{exercise}

\chapter{Filters, Onsets, and Beats \emph{(outline)}}\label{ch:filters}
\chapterlede{How the project isolates frequency bands, detects the instant a
drum is struck, and infers a tempo grid from a noisy sequence of onsets.}

\begin{incode}
Status: outline. To be written against \texttt{feature\_extractor.py} (onset
envelopes, beat tracking) and the legacy Butterworth drum-band split that
DrumSep replaced.
\end{incode}

Planned contents:
\begin{itemize}
  \item \textbf{LTI systems and convolution.} Impulse response, frequency
  response as the DFT of the kernel; why filtering is multiplication in
  frequency (convolution theorem, with proof).
  \item \textbf{Butterworth filters.} Maximally-flat magnitude response
  $\abs{H(\omega)}^2 = 1/(1 + (\omega/\omega_c)^{2n})$; biquad
  implementation; phase distortion and why the legacy kick/snare/hat split
  smeared transients.
  \item \textbf{Onset detection.} Spectral flux
  $\sum_k \max(0, \abs{X[m,k]} - \abs{X[m-1,k]})$; logarithmic compression;
  adaptive thresholding (moving median) --- and the scene-mode variant: the
  per-piece peak-normalized onset gate in \texttt{scene.js} that lets a quiet
  hi-hat register next to a loud kick.
  \item \textbf{Beat tracking as dynamic programming.} Tempo estimation via
  the autocorrelation of the onset envelope; the Ellis DP recurrence
  $D(t) = O(t) + \max_\tau \big(D(t - \tau) - \lambda
  \log^2(\tau/\tau_0)\big)$; backtracking; why breakdowns and rubato break
  it.
  \item \textbf{Sections and phrases.} Self-similarity matrices over chroma
  and timbre features; novelty kernels; the checkerboard detector.
\end{itemize}


% ================= PART II =================
\part{Learning}\label{part:learning}
\chapter{Attention and Transformers \emph{(outline)}}\label{ch:attention}
\chapterlede{The workhorse architecture behind every neural network in this
project: MERT, CLAP's text tower, Demucs' bottleneck, and BS-RoFormer are all
transformers. This chapter builds the mechanism once so later chapters can
use it freely.}

\begin{incode}
Status: outline. The definitions here are prerequisites for
Chapters~\ref{ch:selfsup}--\ref{ch:contrastive}.
\end{incode}

Planned contents:
\begin{itemize}
  \item \textbf{Scaled dot-product attention.}
  $\operatorname{Attn}(Q,K,V) = \softmax(QK^\T/\sqrt{d_k})\,V$; attention as
  a data-dependent convex combination; the $\sqrt{d_k}$ scaling justified via
  the variance of dot products of random vectors (with proof).
  \item \textbf{Multi-head attention} as subspace projections; parameter
  counting; why heads specialize.
  \item \textbf{Positional information.} Sinusoidal encodings; relative
  position; \emph{rotary embeddings} (RoPE) --- rotation matrices acting on
  pairs of coordinates, and the theorem that RoPE makes attention scores a
  function of relative offset only (this is the ``Ro'' in BS-RoFormer).
  \item \textbf{The full block.} Residual streams, LayerNorm (with the
  gradient argument for pre-norm), feed-forward expansion; $O(T^2)$ cost in
  sequence length and what that means for $75$\,fps audio frames.
  \item \textbf{Softmax and temperature.} Properties used throughout:
  invariance to additive shifts, the Jacobian, low/high temperature limits.
\end{itemize}

\chapter{Self-Supervised Audio Representations: MERT \emph{(outline)}}\label{ch:selfsup}
\chapterlede{How a network learns what music sounds like with no labels at
all: mask part of the signal, force the model to predict it, and the
representations needed to succeed turn out to be the ones we wanted.}

\begin{incode}
Status: outline. To be written against \texttt{embeddings.py}
(\texttt{m-a-p/MERT-v1-95M}, frame embeddings at $\sim$75 fps, 13 hidden
layers averaged, mean-pooled per moment).
\end{incode}

Planned contents:
\begin{itemize}
  \item \textbf{The masked-prediction template.} BERT-style objectives for
  continuous signals; why predicting raw waveforms fails and targets must be
  discretized or abstracted.
  \item \textbf{Discrete targets via vector quantization.} Codebooks,
  the RVQ (residual VQ) construction used by neural codecs; commitment
  losses; MERT's acoustic teacher.
  \item \textbf{MERT's musical teacher.} The constant-Q transform (CQT):
  logarithmically-spaced filters aligned with equal temperament
  (Appendix~\ref{ch:temperament}); why adding a CQT reconstruction loss
  injects pitch awareness that codec tokens miss.
  \item \textbf{The convolutional front end.} Strided convolutions as
  learned filterbanks; receptive field and frame-rate arithmetic.
  \item \textbf{Which layer knows what.} Probing studies; why the pipeline
  averages all 13 hidden states rather than taking the last; mean-pooling as
  a bag-of-frames assumption and what temporal information it forfeits.
  \item \textbf{From frames to moments.} The pooling map
  $e = \frac{1}{T}\sum_t h_t$ and its effect on the geometry of
  Chapter~\ref{ch:timbre}.
\end{itemize}

\chapter{Source Separation \emph{(outline)}}\label{ch:separation}
\chapterlede{The signature capability of the whole project: given a stereo
mixture, recover drums, bass, vocals, and the rest. Mathematically: invert a
many-to-one map, using learned priors about what instruments sound like.}

\begin{incode}
Status: outline. To be written against \texttt{ingest.py}: Demucs
(\texttt{htdemucs}, \texttt{htdemucs\_6s}), BS-RoFormer vocals, the MDX23C
DrumSep cascade, and the measured $\sim$17$\times$ vocal-bleed reduction from
the RoFormer-first cascade.
\end{incode}

Planned contents:
\begin{itemize}
  \item \textbf{The problem is ill-posed.} Mixture $m = \sum_s x_s$; the
  solution set is an affine subspace; separation = choosing by prior.
  \item \textbf{Mask-based separation.} Time-frequency masks
  $\hat{X}_s = M_s \odot X_m$; the ideal ratio mask; phase reuse and its
  artifacts; why overlap-add consistency (Exercise 1.4) matters.
  \item \textbf{Waveform models and hybrids.} Demucs' U-Net encoder/decoder
  with a transformer bottleneck operating in both time and spectrogram
  domains; skip connections as multiresolution identity paths.
  \item \textbf{Band-split RoFormer.} Partitioning the spectrogram into
  perceptually-sized bands, per-band tokens, RoPE attention across time and
  band axes (Chapter~\ref{ch:attention}); why this architecture currently
  wins vocals benchmarks.
  \item \textbf{Evaluation.} SDR/SIR/SAR decomposition, in decibels; what a
  $+1$ dB SDR gain sounds like; why ``piano'' is hard (spectral overlap with
  \emph{everything} --- the boundary-flicker documented in the project
  memory).
  \item \textbf{Cascades.} Removing a well-modeled source first
  (vocals) shrinks the residual problem for the rest; conditions under which
  cascading helps vs.\ compounds errors, with the project's drum-kit split
  as the worked example.
\end{itemize}

\chapter{Contrastive Embeddings: CLAP and the Vibe Axes \emph{(outline)}}\label{ch:contrastive}
\chapterlede{One model, two encoders, one shared space: CLAP aligns audio
with language, which lets us name directions in sound-space using English
sentences. The vibe visualizations are driven by dot products against those
directions.}

\begin{incode}
Status: outline. To be written against \texttt{backfill\_clap.py} (audio
tower), \texttt{compute\_vibe.py} (text axes: bright/warm/dense/tense/%
euphoric/vast), and \texttt{microscope.py}'s \texttt{/api/textsearch}.
\end{incode}

Planned contents:
\begin{itemize}
  \item \textbf{The InfoNCE objective.} Batch of $(a_i, t_i)$ pairs;
  \[
  \mathcal{L} = -\frac{1}{B}\sum_i \log
  \frac{\exp(\ip{a_i, t_i}/\tau)}{\sum_j \exp(\ip{a_i, t_j}/\tau)} ;
  \]
  symmetric audio$\to$text and text$\to$audio terms; the learned temperature
  $\tau$; InfoNCE as a lower bound on mutual information (statement, proof
  sketch).
  \item \textbf{Geometry of the shared space.} Why both towers L2-normalize;
  alignment vs.\ uniformity decomposition of the loss; the modality gap
  (audio and text occupy offset cones) and why \emph{differences} of text
  embeddings partially cancel it.
  \item \textbf{Zero-shot classification = nearest text.} Text search over
  moments as a single matrix--vector product (\texttt{emb @ text\_vec}).
  \item \textbf{Semantic axes.} The vibe construction: multi-prompt pole
  averaging (variance reduction), score $= \ip{e, v_+} - \ip{e, v_-}$;
  percentile normalization as a monotone calibration; why axes are
  \emph{named} but not \emph{orthogonal} --- correlated axes (bright/euphoric)
  and what that implies for palette mappings.
  \item \textbf{Failure modes.} Prompt sensitivity; the broken-checkpoint
  incident (\texttt{laion/larger\_clap\_music} collapsed embeddings) as a
  case study in validating learned components empirically.
\end{itemize}


% ================= PART III =================
\part{Geometry of Timbre}\label{part:geometry}
\chapter{Maps of Timbre Space: PCA and UMAP}\label{ch:timbre}
\chapterlede{The galaxy view draws every moment of every song as a star, and
similar-sounding moments land near each other. This chapter contains the
complete mathematics of that map: from $768$-dimensional learned embeddings,
through an optimal linear compression (PCA), to a nonlinear graph embedding
(UMAP) whose axes provably mean nothing --- and whose neighborhoods mean
everything.}

\section{The Data: Moments on a Sphere}

Chapter~\ref{ch:selfsup} describes how MERT turns each $\sim$4-second moment
of audio into a vector $e \in \R^{768}$ by mean-pooling transformer states.
Here we take those vectors as given. The pipeline first normalizes:
\[
x \;=\; \frac{e}{\norm{e}} \;\in\; S^{767},
\]
placing every moment on the unit sphere. This makes similarity purely
\emph{angular}:

\begin{definition}[Cosine similarity and distance]\label{def:cosine}
For nonzero $u, v \in \R^d$,
\[
\operatorname{sim}(u,v) = \frac{\ip{u,v}}{\norm{u}\norm{v}},
\qquad
d_{\cos}(u,v) = 1 - \operatorname{sim}(u,v) \;\in\; [0, 2].
\]
For unit vectors, $\operatorname{sim}(u,v) = \ip{u,v}$ and
$\norm{u - v}^2 = 2\, d_{\cos}(u,v)$: cosine distance and squared Euclidean
distance are the same thing on the sphere, up to a factor of $2$.
\end{definition}

\begin{remark}
$d_{\cos}$ is not a metric (it fails the triangle inequality in general), but
the angle $\arccos(\operatorname{sim})$ is. In practice every algorithm below
only ever \emph{ranks} distances, for which the distinction is harmless.
\end{remark}

\begin{incode}
\texttt{compute\_galaxy.py} stacks the \texttt{emb\_mix} arrays of every song
into $X \in \R^{n \times 768}$ and normalizes rows
(\texttt{X /= np.linalg.norm(X, axis=1, keepdims=True)}). With the current
local library, $n = 664$; the PC batch library will push $n$ into the tens of
thousands, which is where every guarantee in this chapter gets sharper.
\end{incode}

\section{Principal Component Analysis}

The first reduction $768 \to 64$ is linear, classical, and optimal in two
senses that turn out to be the same sense.

Let $X \in \R^{n \times d}$ have centered columns (subtract the mean row
$\bar{x}$ first), and define the sample covariance
\[
C = \frac{1}{n-1} X^\T X \;\in\; \R^{d \times d}.
\]
$C$ is symmetric positive semidefinite, so by the spectral theorem it has an
orthonormal eigenbasis $v_1, \dots, v_d$ with eigenvalues
$\lambda_1 \ge \lambda_2 \ge \cdots \ge \lambda_d \ge 0$.

\begin{theorem}[Variance maximization]\label{thm:pca}
The direction of maximal projected variance,
\[
\argmax_{\norm{u} = 1} \Var(Xu) = \argmax_{\norm{u} = 1} u^\T C u,
\]
is $v_1$, and the maximum equals $\lambda_1$. More generally, the
$k$-dimensional subspace maximizing total projected variance is
$\operatorname{span}(v_1, \dots, v_k)$, capturing variance
$\sum_{i=1}^{k} \lambda_i$.
\end{theorem}

\begin{proof}
Expand $u$ in the eigenbasis: $u = \sum_i \alpha_i v_i$ with
$\sum_i \alpha_i^2 = 1$. Then
\[
u^\T C u = \sum_{i} \lambda_i \alpha_i^2
\;\le\; \lambda_1 \sum_i \alpha_i^2 = \lambda_1,
\]
with equality iff $\alpha_i = 0$ whenever $\lambda_i < \lambda_1$. For the
subspace statement, apply the same argument inductively on the orthogonal
complement of the span chosen so far (the Courant--Fischer characterization).
\end{proof}

\begin{theorem}[Best low-rank approximation, Eckart--Young]\label{thm:eckart}
Write the (thin) singular value decomposition $X = U \Sigma V^\T$ with
singular values $\sigma_1 \ge \cdots \ge \sigma_d$. Among all matrices $M$ of
rank at most $k$,
\[
\norm{X - M}_F^2 \;\ge\; \sum_{i > k} \sigma_i^2,
\]
with equality for $M = U \Sigma_k V^\T$ ($\Sigma_k$ keeps the top $k$
singular values). Equivalently: projecting the rows of $X$ onto
$\operatorname{span}(v_1,\dots,v_k)$ is the rank-$k$ linear map minimizing
total squared reconstruction error.
\end{theorem}

\begin{remark}
Theorems~\ref{thm:pca} and \ref{thm:eckart} are the same fact seen from two
sides, connected by $C = \frac{1}{n-1} V \Sigma^2 V^\T$: the covariance
eigenvectors are the right singular vectors, and
$\lambda_i = \sigma_i^2/(n-1)$. ``Keep the most variance'' and ``lose the
least reconstruction'' pick the same subspace.
\end{remark}

\begin{incode}
\texttt{PCA(n\_components=64).fit\_transform(X)} in
\texttt{compute\_galaxy.py}. Two engineering motives, neither about
visualization per se: (i) UMAP's neighbor search runs $\sim$12$\times$
faster in $64$ dimensions than in $768$; (ii) the trailing dimensions of
mean-pooled transformer embeddings are mostly noise, and
Theorem~\ref{thm:eckart} says truncation is the \emph{least} destructive way
to drop them. Note that with $n = 664 < 768$ the data matrix cannot have rank
above $663$ anyway; PCA here is denoising, not compression.
\end{incode}

\section{UMAP I: The Graph}

UMAP\footnote{McInnes, Healy, Melville, \emph{UMAP: Uniform Manifold
Approximation and Projection for Dimension Reduction} (2018). These notes
present the algorithm as implemented, with the category-theoretic framing
omitted.} operates in two acts. Act one converts the point cloud
$\{x_1, \dots, x_n\} \subset \R^{64}$ into a weighted graph encoding local
structure; act two lays that graph out in the plane. The high-dimensional
coordinates never appear again after act one --- everything the map knows is
in the graph.

\begin{definition}[$k$-NN graph]\label{def:knn}
Fix $k \in \N$ (here $k = 15$) and a metric $d$ (here $d_{\cos}$). For each
$i$ let $N(i)$ be the set of its $k$ nearest neighbors. The directed graph
has an edge $i \to j$ for each $j \in N(i)$.
\end{definition}

Exact $k$-NN costs $\Theta(n^2 d)$; the implementation uses NN-descent, a
randomized refinement scheme, for near-linear approximate search.

The heart of UMAP is what it does with the raw distances next. Two constants
are computed per point:

\begin{definition}[Local calibration]\label{def:rhosigma}
For each point $i$, set $\rho_i = \min_{j \in N(i)} d(x_i, x_j)$ (distance to
the nearest neighbor), and let $\sigma_i > 0$ solve
\begin{equation}\label{eq:sigma}
\sum_{j \in N(i)} \exp\!\left(- \frac{\max(0,\, d(x_i, x_j) - \rho_i)}
{\sigma_i}\right) \;=\; \log_2 k .
\end{equation}
The directed membership strength of edge $i \to j$ is then
\[
w_{i \to j} = \exp\!\left(- \frac{\max(0,\, d(x_i, x_j) - \rho_i)}
{\sigma_i}\right).
\]
\end{definition}

\begin{lemma}[The calibration is well-posed]\label{lem:sigma}
Suppose not all neighbors of $i$ are at distance exactly $\rho_i$ (generic
position). Then Equation~\eqref{eq:sigma} has a unique solution
$\sigma_i \in (0, \infty)$.
\end{lemma}

\begin{proof}
Let $g(\sigma)$ denote the left side, a finite sum of terms
$\exp(-c_j / \sigma)$ with constants $c_j = \max(0, d_{ij} - \rho_i) \ge 0$,
at least one $c_j = 0$ (the nearest neighbor) and, by assumption, at least
one $c_j > 0$. Each term is continuous and nondecreasing in $\sigma$; terms
with $c_j > 0$ are strictly increasing. Hence $g$ is continuous and strictly
increasing, with
\[
\lim_{\sigma \to 0^+} g(\sigma) = \#\{j : c_j = 0\} \ge 1,
\qquad
\lim_{\sigma \to \infty} g(\sigma) = k .
\]
Since $1 \le \log_2 k < k$ for $k \ge 2$ (here $\log_2 15 \approx 3.91$),
the intermediate value theorem gives existence whenever
$\#\{j : c_j = 0\} \le \log_2 k$, and strict monotonicity gives uniqueness.
(The implementation solves \eqref{eq:sigma} by binary search.)
\end{proof}

\begin{remark}[Why calibrate?]\label{rem:uniform}
Equation~\eqref{eq:sigma} forces every point to distribute the same total
membership mass ($\log_2 k$) over its neighborhood. In a dense region,
$\sigma_i$ comes out small (distances are shrunk aggressively); in a sparse
region, large. This is the ``Uniform'' in UMAP: \emph{assume} the data is
uniformly distributed on some manifold, and absorb observed density
variation into a point-dependent rescaling of distance. The $\rho_i$ shift
additionally guarantees every point is connected with weight $1$ to its
nearest neighbor, so no point is orphaned. The cost of this assumption:
UMAP largely erases density information --- a tight cluster and a loose
cluster of the same membership look alike on the map. Cluster \emph{size} on
the galaxy is not evidence of anything.
\end{remark}

Directed weights disagree ($w_{i \to j} \neq w_{j \to i}$), so they are
merged:

\begin{definition}[Fuzzy union]\label{def:fuzzyunion}
The symmetric weight of the pair $\{i, j\}$ is
\[
\mu_{ij} \;=\; w_{i \to j} + w_{j \to i} - w_{i \to j}\, w_{j \to i}.
\]
\end{definition}

\begin{remark}
If the two directed weights are read as independent probabilities that ``$i$
accepts $j$ as a neighbor'' and vice versa, $\mu_{ij}$ is exactly
$\Pr(\text{at least one accepts})$ --- the probabilistic OR
($t$-conorm). One-sided friendships survive; that keeps bridges between
regions alive.
\end{remark}

\section{UMAP II: The Layout}

Act two seeks low-dimensional points $y_1, \dots, y_n \in \R^2$ whose
\emph{own} fuzzy graph resembles $\mu$. The low-dimensional membership is
defined by a smooth rational kernel:

\begin{definition}[Low-dimensional kernel]\label{def:nukernel}
Fix $a, b > 0$. For $y_i, y_j \in \R^2$,
\[
\nu_{ij} \;=\; \Big(1 + a\,\norm{y_i - y_j}^{2b}\Big)^{-1}.
\]
The constants $a, b$ are fit by least squares so that $\nu$ approximates
\[
\Psi(t) =
\begin{cases}
1 & t \le \texttt{min\_dist}, \\
\exp\!\big(-(t - \texttt{min\_dist})\big) & t > \texttt{min\_dist},
\end{cases}
\]
i.e.\ ``full membership within \texttt{min\_dist}, exponential decay
beyond.''
\end{definition}

\texttt{min\_dist} is therefore a pure \emph{rendering} parameter: it sets
how tightly the layout may pack points that the graph considers identical,
trading cluster crispness against within-cluster detail.

\begin{definition}[Objective]\label{def:umapce}
The layout minimizes the fuzzy cross-entropy
\[
\mathcal{C}(y) \;=\; \sum_{i < j}
\underbrace{\mu_{ij} \log \frac{\mu_{ij}}{\nu_{ij}}}_{\text{attraction}}
\;+\;
\underbrace{(1 - \mu_{ij}) \log \frac{1 - \mu_{ij}}{1 - \nu_{ij}}}_{\text{repulsion}} .
\]
\end{definition}

The names are earned by the gradients: the first term is minimized by making
$\nu_{ij}$ large (pull connected points together), the second by making
$\nu_{ij}$ small (push unconnected points apart). The final map is the
equilibrium of $n(n-1)/2$ simultaneous tugs-of-war.

\begin{proposition}[Attractive gradient]\label{prop:grad}
For the attractive term of a single pair, with
$r = \norm{y_i - y_j}$,
\[
\frac{\partial}{\partial y_i}
\Big[ -\mu_{ij} \log \nu_{ij} \Big]
= \mu_{ij}\,
\frac{2ab\, r^{2(b-1)}}{1 + a r^{2b}}\,
\big(y_i - y_j\big),
\]
a spring whose stiffness decays as points approach --- unlike a Hookean
spring, distant neighbors are pulled hard while satisfied ones are left
alone.
\end{proposition}

\begin{proof}
$-\log \nu_{ij} = \log(1 + a r^{2b})$; apply the chain rule with
$\partial r^2 / \partial y_i = 2(y_i - y_j)$ and simplify.
\end{proof}

Minimization proceeds by stochastic gradient descent over \emph{edges}: each
epoch samples edges proportionally to $\mu_{ij}$, applies
Proposition~\ref{prop:grad} to the endpoints, and for each such update also
samples a handful of random non-neighbors ($\approx 5$) to apply the
repulsive gradient to --- \emph{negative sampling}, the standard trick for
avoiding the $\Theta(n^2)$ full repulsion sum. The learning rate anneals
linearly to zero over \texttt{n\_epochs} (the local run used $500$).

Initialization is not random: the points start at the \emph{spectral
embedding} of the graph --- coordinates from the bottom nontrivial
eigenvectors of the symmetric normalized graph Laplacian
$L = I - D^{-1/2} M D^{-1/2}$ (where $M = (\mu_{ij})$ and $D$ is its degree
matrix). This is exactly the Laplacian-eigenmaps solution, a convex
relaxation of the same ``neighbors stay close'' demand, and it starts SGD in
a basin that already respects the coarse cluster structure.

\section{What the Map Does and Does Not Mean}

\begin{proposition}[Isometry invariance]\label{prop:invariance}
Let $T : \R^2 \to \R^2$ be any Euclidean isometry (rotation, reflection,
translation, or a composition). Then
$\mathcal{C}(T y_1, \dots, T y_n) = \mathcal{C}(y_1, \dots, y_n)$.
\end{proposition}

\begin{proof}
$\mathcal{C}$ depends on the layout only through the pairwise distances
$\norm{y_i - y_j}$ (via $\nu_{ij}$), and isometries preserve pairwise
distances by definition.
\end{proof}

\begin{corollary}[The axes are meaningless]\label{cor:axes}
The set of minimizers of $\mathcal{C}$ is closed under all Euclidean
isometries. Which representative the algorithm returns is decided by the
spectral initialization and the randomness of SGD --- not by the data. In
particular, no direction on the galaxy map (up, down, left, right) admits
any interpretation, and separate runs on the same library may return maps
differing by rotation, reflection, or more.
\end{corollary}

What survives, and why one may trust it: $\mu_{ij}$ is built \emph{only}
from each point's $k$ nearest neighbors, and the objective's attraction acts
only on those pairs. Local neighborhoods are the signal; everything else ---
inter-cluster distances, cluster areas (Remark~\ref{rem:uniform}), global
shape --- is a side effect of repulsion and should be read with suspicion.

\begin{incode}
\texttt{compute\_galaxy.py} runs
\texttt{umap.UMAP(n\_components=2, n\_neighbors=15, min\_dist=0.05,
metric="cosine")} for positions. Post-processing: positions are centered and
divided by the $99.5$th percentile of coordinate magnitude --- a robust
rescale to roughly $[-1,1]^2$ that one outlier moment cannot hijack.
The client (\texttt{app.js}, \texttt{Galaxy}) and the drift view
(\texttt{vibescenes.js}) treat the coordinates as opaque: everything they
draw --- hover neighborhoods, the song path, the flight camera --- uses only
relative position, as Corollary~\ref{cor:axes} demands.
\end{incode}

\section{Color as a Second Opinion}

The star colors are not styling; they are a second, independent instance of
the same mathematics. The pipeline runs UMAP \emph{again} on the same
$64$-dimensional data with \texttt{n\_components=3} and a looser
\texttt{min\_dist} of $0.3$, yielding $z_i \in \R^3$, then maps robustly to
RGB:
\[
\text{rgb}_i \;=\; 60 + 195 \cdot
\operatorname{clip}\!\left(
\frac{z_i - p_2}{p_{98} - p_2},\; 0,\; 1\right),
\]
where $p_2, p_{98}$ are per-channel percentiles (the affine map lands in
$[60, 255]^3$ so no star is too dark against the background).

Because the 3-D run shares the input graph-building mathematics but not the
2-D run's optimization trajectory, \emph{hue similarity is a
quasi-independent witness of timbral similarity}. When two nearby stars also
share a color, two different embeddings agreed; when a star's color clashes
with its neighborhood, the 2-D layout likely forced a compromise
(Corollary~\ref{cor:axes} again: the plane is small).

\section*{Exercises}

\begin{exercise}
Prove the identity $\norm{u - v}^2 = 2\,d_{\cos}(u,v)$ for unit vectors, and
exhibit three unit vectors in $\R^2$ violating the triangle inequality for
$d_{\cos}$.
\end{exercise}

\begin{exercise}
Prove the subspace half of Theorem~\ref{thm:pca}: among all $k$-dimensional
subspaces $S$, total projected variance is maximized by the span of the top
$k$ eigenvectors. (Induct using the variational characterization on the
orthogonal complement.)
\end{exercise}

\begin{exercise}
In Lemma~\ref{lem:sigma}, suppose point $i$ has $t$ neighbors tied exactly at
distance $\rho_i$ with $t > \log_2 k$. What goes wrong, and what does the
condition mean geometrically? (The implementation clamps $\sigma_i$ to a
minimum value in this case.)
\end{exercise}

\begin{exercise}
Verify Proposition~\ref{prop:grad}. Then derive the repulsive gradient
$\partial\big[-(1-\mu_{ij})\log(1 - \nu_{ij})\big] / \partial y_i$ and show
it diverges as $r \to 0$ when $b < 1$. The implementation adds a small
$\varepsilon$ to $r^2$; explain what failure this prevents during SGD.
\end{exercise}

\begin{exercise}
Prove that $\mu_{ij} = w_{i\to j} + w_{j \to i} - w_{i \to j} w_{j \to i}$
satisfies $\max(w_{i \to j}, w_{j \to i}) \le \mu_{ij} \le
\min(1,\, w_{i \to j} + w_{j \to i})$, and characterize when each bound is
attained.
\end{exercise}

\begin{exercise}[computational]
Run \texttt{compute\_galaxy.py} twice with different random seeds and overlay
the two maps. Find the best-fitting rotation+reflection between them by
solving the orthogonal Procrustes problem
$\min_{Q^\T Q = I} \norm{Y_1 Q - Y_2}_F$ (solution: $Q = UV^\T$ from the SVD
of $Y_1^\T Y_2$). How much residual disagreement survives after aligning?
Which regions of the map are stable, and are they the dense ones?
\end{exercise}

\begin{exercise}[computational]
Set \texttt{n\_neighbors=5} and rebuild the galaxy. Describe what happens to
(i) song threads, (ii) the number of visual clusters, (iii) bridges between
genres. Explain each observation via Remark~\ref{rem:uniform} and the role of
$k$ in Definition~\ref{def:knn}.
\end{exercise}

\chapter{Similarity Search and the Moment Index \emph{(outline)}}\label{ch:similarity}
\chapterlede{``Find me moments that sound like this one'' --- across every
song in the library, in milliseconds, along a chosen facet (whole mix, just
the drums, just the vocals).}

\begin{incode}
Status: outline. To be written against \texttt{moment\_index.py},
\texttt{moments.py}, \texttt{descriptors.py}, \texttt{interactions.py}, and
the \texttt{/api/similar} endpoint.
\end{incode}

Planned contents:
\begin{itemize}
  \item \textbf{Exact top-$k$ retrieval.} Cosine scores as one matrix--vector
  product; complexity and memory layout; when brute force beats an index
  (spoiler: at this library's scale, always).
  \item \textbf{Faceted similarity.} Per-stem embedding blocks
  (\texttt{emb\_drums}, \texttt{emb\_bass}, \ldots) as projections of the
  product space; facet weighting as a diagonal metric.
  \item \textbf{Interpretable descriptors.} The 18-dimensional
  descriptor vector (RMS, onset density, spectral centroid, active
  fraction per stem); z-scoring; mixing learned and hand-crafted metrics.
  \item \textbf{Diversity constraints.} The per-song result cap as a greedy
  approximation to maximal marginal relevance.
  \item \textbf{Approximate methods} (for the future large library):
  inverted-file indices, HNSW graphs, product quantization --- what each
  trades away.
\end{itemize}


% ================= PART IV =================
\part{Rendering}\label{part:rendering}
\chapter{The Mathematics of Rendering \emph{(outline)}}\label{ch:shaders}
\chapterlede{Everything the visualizer draws is mathematics executed sixty
times a second: procedural noise, divergence-free flow fields, radial
symmetry groups, perspective projection, and color spaces.}

\begin{incode}
Status: outline. To be written against \texttt{scene.js} (mandala shader,
stage projection) and \texttt{vibescenes.js} (liquid advection, weather FBM,
drift camera).
\end{incode}

Planned contents:
\begin{itemize}
  \item \textbf{Hash, value noise, and fBm.} Lattice noise with smoothstep
  interpolation ($C^1$ but not $C^2$ --- and where that artifact is visible);
  fractional Brownian motion $\sum_k a^k\, n(2^k p)$, its spectrum, and the
  self-similarity parameter.
  \item \textbf{Curl noise.} In 2-D, the stream-function construction
  $v = (\partial_y \psi, -\partial_x \psi)$; \emph{theorem:}
  $\nabla \cdot v = 0$ (incompressible flow --- why the liquid view folds
  instead of draining); finite-difference curl and its step-size tradeoff.
  \item \textbf{Semi-Lagrangian advection.} The backward-lookup scheme
  $D_{t+1}(p) = D_t(p - v\,\Delta t)$; unconditional stability; numerical
  diffusion as implicit blur, exploited deliberately as dye decay.
  \item \textbf{Kaleidoscopes as group actions.} The dihedral group
  $D_{12}$ acting on the plane; folding $\theta$ into a fundamental domain
  via $\abs{\theta - s\lfloor \theta/s \rceil}$; why chroma's 12 pitch
  classes make this the \emph{honest} symmetry group for harmony.
  \item \textbf{Perspective projection.} The stage view's camera
  (\texttt{project()} in \texttt{scene.js}): yaw/pitch rotations as matrix
  products, similar triangles for the focal divide, painter's-algorithm
  depth sorting.
  \item \textbf{Color.} HSL/HSV as cylindrical re-parameterizations of RGB;
  hue arithmetic on the circle; why the vibe palette interpolates hue rather
  than RGB (avoiding desaturated midpoints); gamma and additive blending in
  the drift view.
  \item \textbf{Feedback systems.} Ping-pong framebuffers as discrete
  dynamical systems $D_{t+1} = f(D_t) + \text{input}_t$; contraction (decay
  $< 1$) as the guarantee the liquid never blows up.
\end{itemize}


% ================= appendix =================
\appendix
\chapter{Equal Temperament and Chroma}\label{ch:temperament}
\chapterlede{The small amount of music theory the codebase depends on,
stated as mathematics.}

\section{Pitch is Logarithmic}

The ear judges two pitch \emph{intervals} equal when the frequency
\emph{ratios} are equal: the step from $110$ to $220$ Hz sounds like the
step from $220$ to $440$ Hz. Perceived pitch is therefore a logarithm of
frequency, and Western tuning discretizes that logarithm.

\begin{definition}[12-tone equal temperament]\label{def:tet}
Fix the reference $A4 = 440$ Hz. The \emph{semitone} is the ratio
$2^{1/12}$, and the pitch number (MIDI number) of frequency $f > 0$ is
\[
m(f) \;=\; 69 + 12 \log_2\!\frac{f}{440},
\]
so that $m(440) = 69$ and each octave (doubling) spans exactly $12$ units.
The nearest note is $\lfloor m(f) \rceil$, and the error
$m(f) - \lfloor m(f) \rceil$, in units of $1/100$ of a semitone, is measured
in \emph{cents}.
\end{definition}

\begin{remark}
$2^{1/12}$ is irrational, so no interval except the octave is acoustically
pure: the equal-tempered fifth is $2^{7/12} \approx 1.4983$, not $3/2$. Equal
temperament trades purity for symmetry --- every key sounds the same --- and
it is exactly that symmetry the definitions below exploit.
\end{remark}

\begin{definition}[Pitch class and chroma]\label{def:chroma}
The \emph{pitch class} of frequency $f$ is
$c(f) = \lfloor m(f) \rceil \bmod 12 \in \{0, \dots, 11\}$ ($0 = $ C,
$1 = $ C$\sharp$, \dots, $11 = $ B). A \emph{chroma vector} of a spectrum
$\abs{X[k]}$ is the 12-vector that folds all octaves onto pitch classes:
\[
\text{ch}[c] \;=\; \sum_{k \,:\, c(f_k) = c} \abs{X[k]},
\qquad c = 0, \dots, 11,
\]
summing over bins $k$ with frequencies $f_k$ in a chosen band.
\end{definition}

Chroma deliberately forgets octave and timbre, keeping only \emph{harmony}
--- which notes are sounding. Two voicings of the same chord, octaves apart
on different instruments, give nearly the same chroma vector.

\begin{incode}
Definition~\ref{def:tet} appears verbatim wherever the code names a note:
\texttt{scene.js} computes \texttt{midi = Math.round(69 + 12 *
Math.log2(hz / 440))} for the vocal and bass note labels, with
\texttt{NOTE\_NAMES[((midi \% 12) + 12) \% 12]} implementing $c(f)$ (the
double-mod guards JavaScript's signed \texttt{\%}). Definition~\ref{def:chroma}
is \texttt{computeChroma()}, folding analyser bins between $55$ and
$2000$ Hz; the twelve values drive the mandala's twelve petals, making that
visualization literally a plot of the chroma vector in polar coordinates.
The per-bin nearest-note rounding is exact for high bins but aliases for low
ones (Chapter~\ref{ch:sound}, Exercise 1.6) --- one reason chroma starts at
$55$ Hz rather than lower.
\end{incode}

\section*{Exercises}

\begin{exercise}
Show that the map $f \mapsto c(f)$ satisfies $c(2f) = c(f)$, and that scaling
$f$ by the equal-tempered fifth $2^{7/12}$ advances the pitch class by $7
\pmod{12}$. Since $\gcd(7, 12) = 1$, conclude that stacking fifths visits all
$12$ pitch classes --- the circle of fifths as a statement about the group
$\Z/12\Z$.
\end{exercise}

\begin{exercise}
The bass string $E1$ has fundamental $41.2$ Hz. Its second and third
harmonics are $82.4$ and $123.6$ Hz. To which pitch classes do the first
five harmonics round, and how many cents does each miss by? (This is why
chroma computed on a bass-heavy band is polluted by harmonics --- and why
\texttt{computeChroma} prefers harmonic stems.)
\end{exercise}


\end{document}
